% ============================================================
%  RÉSUMÉS (4e de couverture)
% ============================================================
\chapter*{}
\thispagestyle{empty}

\vspace*{\fill}

\begin{center}
{\large\bfseries\color{marron}
Digitalisation de l'expérience client pour l'huile d'olive extra vierge\\
via un configurateur interactif à visualisation 3D
}\\[0.3cm]
{\normalsize \NomEtudiant}
\end{center}

\vspace{1cm}

% ---- Résumé Arabe ----
\begin{flushright}
\textbf{:ملخص}\\[0.3cm]
\begin{minipage}{0.95\textwidth}
\begin{flushright}
تمت كتابة هذه الرسالة في إطار مشروع نهاية الدراسة للحصول على شهادة الإجازة في الإعلامية والوسائط المتعددة. الهدف من هذا العمل هو تطوير منصة ويب تفاعلية مخصصة لتثمين زيت الزيتون التونسي البكر الممتاز لعلامة "فندري". تمت النمذجة المفاهيمية باستخدام لغة UML. أما التنفيذ فقد اعتمد على React وTypeScript في الواجهة الأمامية، وNode.js وExpress.js في الواجهة الخلفية، فضلاً عن قاعدة بيانات MongoDB. كما تم دمج نماذج ثلاثية الأبعاد واقعية للمنتجات، تم إنشاؤها باستخدام برنامج Blender، في واجهة مُهيِّج تفاعلي يتيح للمستخدم تخصيص قنينة زيت الزيتون واختيار التغليف والحجم والتصميم.
\end{flushright}
\end{minipage}\\[0.3cm]
\textbf{المفاتيح:} تجارة إلكترونية، مُهيِّج تفاعلي، نمذجة ثلاثية الأبعاد، TypeScript، UML، MongoDB، Node.js، React
\end{flushright}

\vspace{0.8cm}
\textcolor{or}{\rule{\textwidth}{0.4pt}}
\vspace{0.8cm}

% ---- Résumé Français ----
\noindent\textbf{Résumé :} Ce mémoire a été rédigé dans le cadre du projet de fin d'études pour l'obtention du diplôme de licence en informatique et multimédia. L'objectif de ce travail est de développer une plateforme web interactive dédiée à la valorisation de l'huile d'olive extra vierge tunisienne de la marque Domaine FENDRI. La modélisation conceptuelle de cette application a été réalisée avec le langage de conception et de modélisation UML. La réalisation du projet a été développée avec React et TypeScript dans le frontend, Node.js et Express.js dans le backend, et MongoDB comme système de gestion de base de données. Des modèles 3D photoréalistes des produits, conçus sous Blender, ont été intégrés dans un configurateur interactif permettant à l'utilisateur de personnaliser sa bouteille d'huile d'olive selon le contenant, l'étiquette, l'emballage et la quantité.

\noindent\textbf{Mots-clés :} E-commerce, configurateur interactif, visualisation 3D, React, Node.js, MongoDB, UML, TypeScript.

\vspace{0.5cm}
\textcolor{or}{\rule{\textwidth}{0.4pt}}
\vspace{0.5cm}

% ---- Abstract Anglais ----
\noindent\textbf{Abstract:} This thesis was written as part of the end-of-studies project for obtaining the bachelor's degree in computer science and multimedia. The objective of this work is to develop an interactive web platform dedicated to the digital valorisation of Tunisian extra virgin olive oil for the Domaine FENDRI brand. The conceptual modeling of this application was carried out using the UML design and modeling language. The project was developed using React and TypeScript for the frontend, Node.js and Express.js for the backend, and MongoDB as the database management system. Photorealistic 3D models of the products, created with Blender, were integrated into an interactive configurator allowing the user to customize their olive oil bottle by selecting the container, label, packaging, and quantity.

\noindent\textbf{Key-words:} E-commerce, interactive configurator, 3D visualization, React, Node.js, MongoDB, UML, TypeScript.

\vspace{\fill}

\begin{center}
\textcolor{gris}{\small
Pôle technologique route Tunis Km 10 ; B.P. : 242 SFAX 3021\\
Tél. : 216 74 86 22 33 ; Fax : 216 74 86 24 32\\
\url{www.isimsf.rnu.tn}
}
\end{center}
